\chapter*{Ilustracije}
\addcontentsline{toc}{chapter}{Ilustracije}

\section*{Glava 1 — Poziv na Hiperion}

\textbf{Fajl:} \texttt{images/01-poziv-na-hiperion.jpg}

\textbf{Mesto:} Hiperion.

\textbf{Scena:} Konzul u svom izolovanom domu iznad močvare, okružen ogromnim drvećem dok se približava oluja.

\textbf{Likovi:} Konzul.

\textbf{Funkcija ilustracije:} Uspostavlja osećaj samoće, predosećaja i trenutka neposredno pre poziva koji pokreće radnju.

\section*{Glava 2 — Poslednje hodočašće}

\textbf{Fajl:} \texttt{images/02-poslednje-hodocasce.jpg}

\textbf{Mesto:} Igdrasil, tokom puta ka Hiperionu.

\textbf{Scena:} Sedam hodočasnika prvi put okupljenih kao grupa.

\textbf{Likovi:} Konzul, Lenar Hojt, Fedman Kasad, Martin Silenus, Sol Vejntraub sa Rejčel, Bron Lamija i Het Mastin.

\textbf{Funkcija ilustracije:} Predstavlja grupu pre nego što počnu da dele svoje priče i naglašava njihove različite karaktere.

\section*{Glava 3 — Sveštenik}

\textbf{Fajl:} \texttt{images/03-svestenik.jpg}

\textbf{Mesto:} Pukotina na Hiperionu.

\textbf{Scena:} Pol Dire među Bikurom, dok pokušava da razume njihovu religiju i tajnu kruciforme.

\textbf{Likovi:} Pol Dire i Bikure.

\textbf{Funkcija ilustracije:} Nagoveštava religioznu misteriju, izolaciju i nelagodu pre punog otkrivanja prirode Bikura.

\section*{Glava 4 — Vojnik}

\textbf{Fajl:} \texttt{images/04-vojnik.jpg}

\textbf{Mesto:} Grad Pesnika na Hiperionu.

\textbf{Scena:} Fedman Kasad i Moneta u napuštenom Gradu Pesnika.

\textbf{Likovi:} Fedman Kasad i Moneta.

\textbf{Funkcija ilustracije:} Spaja Kasadovu vojničku prošlost sa njegovom opsesijom Monetom i misterijom Šrajka.

\section*{Glava 5 — Pesnik}

\textbf{Fajl:} \texttt{images/05-pesnik.jpg}

\textbf{Mesto:} Grad Pesnika na Hiperionu.

\textbf{Scena:} Martin Silenus piše \emph{Spev o Hiperionu} među ruševinama napuštene umetničke kolonije.

\textbf{Likovi:} Martin Silenus; Šrajk može biti samo diskretno nagovešten.

\textbf{Funkcija ilustracije:} Predstavlja vezu između umetničkog stvaranja, opsesije i smrti.

\section*{Glava 6 — Naučnik}

\textbf{Fajl:} \texttt{images/06-naucnik.jpg}

\textbf{Mesto:} Barnardov svet.

\textbf{Scena:} Sol Vejntraub sa Rejčel, koja se zbog Merlinove bolesti vraća kroz sopstveno detinjstvo.

\textbf{Likovi:} Sol i Rejčel.

\textbf{Funkcija ilustracije:} Naglašava roditeljsku ljubav, nemoć i gubitak koji čine emotivno jezgro Solove priče.

\section*{Glava 7 — Detektivka}

\textbf{Fajl:} \texttt{images/07-detektivka.jpg}

\textbf{Mesto:} Lusus.

\textbf{Scena:} Bron Lamija i Džoni na početku istrage.

\textbf{Likovi:} Bron Lamija i Džoni.

\textbf{Funkcija ilustracije:} Uvodi noir ton, pitanje identiteta i Džonijevu neobičnu prirodu.

\section*{Glava 8 — Konzul}

\textbf{Fajl:} \texttt{images/08-konzul.jpg}

\textbf{Mesto:} Maui-Kovenant.

\textbf{Scena:} Siri i Merin uz more, u svetu koji će vremenski dug i dolazak Hegemonije promeniti.

\textbf{Likovi:} Siri i Merin Aspik.

\textbf{Funkcija ilustracije:} Predstavlja ljubav, vremensku udaljenost i gubitak Maui-Kovenanta.

\section*{Glava 9 — Dolina Vremenskih grobnica}

\textbf{Fajl:} \texttt{images/09-dolina-vremenskih-grobnica.jpg}

\textbf{Mesto:} Dolina Vremenskih grobnica, Hiperion.

\textbf{Scena:} Hodočasnici silaze zajedno prema Vremenskim grobnicama.

\textbf{Likovi:} Šest odraslih hodočasnika i Rejčel; Het Mastin nije sa grupom.

\textbf{Funkcija ilustracije:} Zatvara prvi deo knjige slikom zajedništva, straha i ulaska u nepoznato.